\scanpage{216}
on a

\[
(z,\varphi_\lambda(z))\in R\iff\varphi_\lambda(z)=z\quad\text{i.e.}\quad2z=\lambda \tag*{(73)}
\]

et \deleted{rien n'empêche qu'il y ait} \inserted{il y a} des solutions \inserted{en $z$} de cette équation pour certaines valeurs du paramètre (\deleted{\unsure{notamment}} \inserted{savoir} pour $\lambda\in2Z$ -- NB on écrit additivement ici la loi de groupe), alors qu'il est possible qu'il n'y en ait pas pour d'autres (dès lors que $2Z\ne Z$) -- i.e. qu'on peut \inserted{néanmoins} choisir une ``origine'' $\varphi_{\lambda_0}$ telle que elle \unsure{si|se} corresponde à une ``vraie'' structure de graphe $A_{\lambda_0}$. L'exemple le plus économique serait \deleted{$\unsure{E}$} $Z=\mathbb Z/2\mathbb Z$ \unsure{,} \ldots

Soit $\Pi$ un sous-$Z$-torseur du $\mathcal G$-torseur des $Z$-jauges du prégraphe $\Gamma=(S,S^\vee,\Gamma)$. Ainsi $\unsure{\Pi|\Gamma}$ est formé de bijections

\[
\varphi:S\mathrel{\overset{\sim}{\rule{2em}{0.35pt}}} S^\vee
\]

satisfaisant

\[
\lambda_{S^\vee}^{\unsure{*|\vee}}\varphi(s)=\deleted{\lambda_S}\varphi(\lambda_S^{-1}(s))
\qquad \lambda_{S^\vee}^{*}\varphi=\varphi\circ\lambda_S^{-1} \tag*{(74)}
\]

i.e. compatibles avec les opérations de $Z$ sur $S$ et $S^\vee$, via l'hom.

\[
\begin{aligned}
&\lambda\longmapsto\lambda^{-1}\\
&Z\longrightarrow Z
\end{aligned} \tag*{(75)}
\]

des groupes d'opérateurs. \deleted{On peut donc considérer que $S^\vee$ est déduit de $S$ par le torseur} Si $\lambda\in Z$, on aura en tordant \deleted{\unsure{par}} des \unsure{rels|deux} \deleted{par le $Z$-torseur $\Pi$, pour l'opération de $Z$ sur $S$ $T_\lambda\varphi=\lambda_{S^\vee}\varphi=\varphi\lambda_S^{-1}$}

% archive sheet number: 215
